\documentclass[12pt,a4paper]{article}
\usepackage[utf8]{inputenc}
\usepackage[brazilian]{babel}
\usepackage{amsmath}
\usepackage{amsfonts}
\usepackage{amssymb}
\usepackage{booktabs}
\usepackage{geometry}
\usepackage[hyphens]{url} 
\usepackage[colorlinks=true, urlcolor=blue, linkcolor=blue]{hyperref}
\usepackage{dirtree} 
\usepackage{xcolor}
\usepackage{graphicx}
\usepackage{enumitem}

\geometry{left=2cm, right=2cm, top=3cm, bottom=2cm}

\hypersetup{
	colorlinks=true,
	urlcolor=blue,
	linkcolor=blue
}

\begin{document}
	
	% Cabeçalho da Avaliação
	\noindent
	\begin{center}
		\LARGE \textbf{Avaliação Prática Unificada (Parte 1)} \\
		\vspace{0.2cm}
		\large \textbf{Construção de Pipeline de Dados com Fundamentação em Ciência de Dados} \\
		\vspace{0.4cm}
		\normalsize \textit{Disciplina: Ciência de Dados} \\
		\textbf{Data de Entrega: 28/06/2026} \\
		\normalsize Formato: Grupos ou Individual
	\end{center}
	
	\vspace{0.2cm}
	\noindent\rule{\textwidth}{0.4pt}
	\vspace{0.5cm}
	
	\section{Objetivo da Avaliação}
	Esta avaliação visa integrar os conhecimentos práticos de manipulação de dados com os fundamentos teóricos de Ciência de Dados (Inferência Causal, Amostragem, Análise Exploratória e Estatística). O objetivo é construir um pipeline reprodutível que não apenas processe os dados, mas que seja projetado com plena consciência sobre a natureza, as limitações estatísticas e a integridade analítica da informação.
	
	\section{Escolha do Tema e Fontes de Dados (Kaggle)}
	Cada equipe deverá escolher \textbf{um} dos temas abaixo. O processo de extração deve obrigatoriamente consumir as fontes indicadas do Kaggle.
	
	\begin{itemize}[leftmargin=*]
		\item \textbf{Opção A: Logística e Segurança Viária} \\
		{\footnotesize \url{https://www.kaggle.com/datasets/rafaelborgesgraunke/traffic-accidents-brazil-pt-br}}
		\item \textbf{Fonte Secundária Auxiliar:} \\
		{\footnotesize \url{https://www.kaggle.com/datasets/liamarguedas/brazil-total-highway-crashes-2010-2023}}
		
		\item \textbf{Opção B: Finanças e E-Commerce} \\
		{\footnotesize \url{https://www.kaggle.com/datasets/olistbr/brazilian-ecommerce}}
		
		\item \textbf{Opção C: Meio Ambiente (Queimadas)} \\
		{\footnotesize \url{https://www.kaggle.com/datasets/gustavomodelli/forest-fires-in-brazil}}
		
		\item \textbf{Opção D: Saúde Pública (COVID-19)} \\
		{\footnotesize \url{https://www.kaggle.com/datasets/gpreda/covid-world-vaccination-progress}}
		
		\item \textbf{Opção E: Educação e Desempenho} \\
		{\footnotesize \url{https://www.kaggle.com/datasets/csafrit2/higher-education-students-performance-evaluation}}
		
		\item \textbf{Opção F: Energia, Clima e Recursos Renováveis} \\
		{\footnotesize \url{https://www.kaggle.com/datasets/pythonafroz/renewable-power-generation-and-weather-conditions}}
	\end{itemize}
	
	\section{Detalhamento das Atividades (Aplicação de Conceitos)}
	
	\subsection{Camada de Ingestão (Amostragem e Viés)}
	Vocês devem construir o script de captura e leitura inicial dos dados brutos (\texttt{extractor.py}).
	\begin{itemize}
		\item \textbf{Desenvolvimento:} Implementar a leitura automatizada dos arquivos \texttt{.csv} originais e isolar falhas de caminhos de arquivos por meio de tratamento de erros. O script deve registrar um log simples informando as dimensões iniciais da matriz de dados (volumetria bruta).
		\item \textbf{Aplicação Teórica (README):} Na documentação do projeto, identifique explicitamente qual seria a \textbf{População-alvo} ideal que os dados do seu tema deveriam representar e qual é a \textbf{Estrutura de Acesso} (\textit{Access Frame}) real oferecida pelo dataset disponível no Kaggle. Discuta criticamente se há riscos de \textbf{Viés de Seleção} na forma como esses dados foram originalmente coletados.
	\end{itemize}
	
	\subsection{Análise Exploratória e Tratamento Estatístico (EDA)}
	Vocês devem estruturar as regras de transformação e higienização das variáveis (\texttt{cleaner.py}).
	\begin{itemize}
		\item \textbf{Desenvolvimento (Sanitização e IQR):} Identificar e tratar falhas de codificação de texto (\textit{encoding}) e padronizar strings categóricas. Implementar algoritmicamente o cálculo estatístico do \textbf{Intervalo Interquartil (IQR)} para as colunas numéricas contínuas mais relevantes, isolando os \textit{outliers} de forma automatizada para evitar distorções estatísticas.
		\item \textbf{Dicionário de Dados:} No \texttt{README.md}, monte uma tabela com o esquema final do dataset. Classifique obrigatoriamente cada variável tratada sob a ótica da modelagem estatística como \textbf{Discreta}, \textbf{Contínua} ou \textbf{Categórica}.
		\item \textbf{Tratamento de Nulos (Viés vs Variância):} Justifique explicitamente no relatório a estratégia adotada (exclusão de registros ou imputação por medidas de tendência central). Explique o impacto teórico dessa decisão no \textbf{Viés} (erros sistemáticos) e na \textbf{Variância} (sensibilidade a novos dados) do seu conjunto final.
		\item \textbf{Consolidação Lógica (Merge):} Realizar o cruzamento das tabelas por chaves comuns (ex: códigos geográficos, IDs). O código deve mitigar disparidades nas chaves para evitar a perda silenciosa de amostras significativas durante o \textit{join}.
	\end{itemize}
	
	\subsection{Análise de Domínio (Inferência Causal e Ceteris Paribus)}
	Esta seção será avaliada exclusivamente pela profundidade conceitual e científica demonstrada na documentação.
	\begin{itemize}
		\item \textbf{Aplicação Teórica (README):} Escolha uma relação hipotética ou lógica entre duas variáveis do tema do seu dataset (ex: \textit{Intensidade de Chuva} e \textit{Volume de Acidentes}; \textit{Renda Familiar} e \textit{Acesso à Educação superior}; \textit{Radiação Solar} e \textit{Eficiência do Painel}). Responda com rigor científico:
		\begin{enumerate}[label=\alph*)]
			\item Explique detalhadamente por que uma forte correlação matemática encontrada entre essas variáveis não é evidência suficiente para comprovar um nexo de causalidade.
			\item Identifique e disserte sobre pelo menos duas possíveis \textbf{Variáveis de Confusão} (\textit{confounders}) omitidas que poderiam enviesar os resultados caso fossem ignoradas.
			\item Descreva conceitualmente como seria o desenho de um cenário ideal sob a ótica do princípio do \textbf{Ceteris Paribus} para isolar e validar adequadamente o efeito causal dessa relação.
		\end{enumerate}
	\end{itemize}
	
	\subsection{Visualização Científica (Integridade Visual)}
	Vocês devem criar uma camada visual dedicada para traduzir graficamente os insights coletados.
	\begin{itemize}
		\item \textbf{Desenvolvimento:} Construa o script \texttt{visualize.py} que utilize a base de dados tratada pelo pipeline para gerar de forma automatizada \textbf{um} dos seguintes gráficos:
		\begin{itemize}
			\item Evolução temporal da métrica principal (Gráfico de Linhas).
			\item Distribuição proporcional de atributos qualitativos (Gráfico de Barras).
			\item Relação e dispersão entre duas variáveis numéricas (Gráfico de Dispersão).
		\end{itemize}
		\item \textbf{Integridade Visual:} O gráfico exportado deve respeitar estritamente os princípios de ética e honestidade visual em Ciência de Dados, contendo eixos devidamente identificados, unidades de medida explícitas, escalas proporcionais coerentes (eixos numéricos iniciando em zero onde aplicável) e ausência total de distorções visuais.
	\end{itemize}
	
	\section{Estrutura do Projeto e Entregáveis}
	O repositório do projeto no \textbf{Git} deve seguir o formato mínimo e organizado especificado abaixo para garantir a reprodutibilidade analítica:
	
	\vspace{0.3cm}
	\begin{minipage}{\textwidth}
		\dirtree{%
			.1 meu-projeto-data-science/.
			.2 src/.
			.3 extract/.
			.4 extractor.py \hfill\begin{minipage}[t]{8.5cm}\textit{Script de ingestão e leitura de arquivos}\end{minipage}.
			.3 transform/.
			.4 cleaner.py \hfill\begin{minipage}[t]{8.5cm}\textit{Script com tratamento estatístico (Nulos e IQR)}\end{minipage}.
			.3 visualize.py \hfill\begin{minipage}[t]{8.5cm}\textit{Script focado em gráficos de integridade visual}\end{minipage}.
			.3 main.py \hfill\begin{minipage}[t]{8.5cm}\textit{Orquestrador central de execução do pipeline}\end{minipage}.
			.2 requirements.txt \hfill\begin{minipage}[t]{8.5cm}\textit{Dependências de bibliotecas de análise}\end{minipage}.
			.2 README.md \hfill\begin{minipage}[t]{8.5cm}\textit{Relatório científico e documentação teórica}\end{minipage}.
		}
	\end{minipage}
	\vspace{0.4cm}
	
	O pipeline de execução deve gerar de forma automatizada o arquivo final limpo na raiz do projeto sob o nome \texttt{dados\_limpos\_final.csv} (ou \texttt{.json}), representando a base consolidada para futuras inferências.
	
	\section{Critérios de Avaliação}
	\begin{center}
		\begin{tabular}{p{11cm}cc}
			\toprule
			\textbf{Critério de Avaliação Científica} & \textbf{Peso} \\
			\midrule
			\textbf{Manipulação Estatística de Dados (EDA):} Qualidade algorítmica no tratamento de nulos, inconsistências de tipos, pareamento de dados e isolamento matemático de outliers via IQR. & 40\% \\
			\addlinespace
			\textbf{Fundamentação de Ciência de Dados:} Rigor conceitual na análise de viés de amostragem, modelagem de variáveis e justificativa das hipóteses de inferência causal e confusão no README. & 40\% \\
			\addlinespace
			\textbf{Visualização e Integridade:} Geração via script automatizado do gráfico escolhido sob os critérios de honestidade estatística e ética visual. & 20\% \\
			\bottomrule
		\end{tabular}
	\end{center}
	
\end{document}